\section{Regroup twenty-four bits into four sextets}
A \emph{sextet} is a group of six bits, whose unsigned value is from zero through 63.
Let the three input bytes be \(b_0,b_1,b_2\). The first byte supplies all six bits of
\(q_0\) and the high two bits of \(q_1\). The next byte finishes \(q_1\) and starts
\(q_2\); the third finishes \(q_2\) and supplies \(q_3\):
\[
\begin{aligned}
q_0&=b_0\gg2,\\
q_1&=((b_0\mathbin{\&}3)\ll4)\mathbin{\mathrm{OR}}(b_1\gg4),\\
q_2&=((b_1\mathbin{\&}15)\ll2)\mathbin{\mathrm{OR}}(b_2\gg6),\\
q_3&=b_2\mathbin{\&}63.
\end{aligned}
\]
The right shifts are logical. Masking the low two or four bits before shifting makes
each contribution occupy only its assigned part of the next sextet. Their bit fields
do not overlap, so OR joins them without a carry.

\begin{figure}[htbp]
\centering
\begin{tikzpicture}[alt={The three input bytes for Man, binary 01001101, 01100001 and 01101110, occupy one 24-bit span. Repartitioning the same span into six-bit groups gives 010011,010110,000101,101110: decimal 19,22,5,46 and characters T,W,F,u.}]
\foreach \x/\bits/\name in {0/01001101/M,4/01100001/a,8/01101110/n} {
 \draw (\x,1.5) rectangle +(4,1);
 \node at (\x+2,2.1) {\texttt{\bits}};
 \node at (\x+2,1.7) {\name};
}
\foreach \x/\bits/\value/\name in {0/010011/19/T,3/010110/22/W,6/000101/5/F,9/101110/46/u} {
 \draw (\x,0) rectangle +(3,1);
 \node at (\x+1.5,.7) {\texttt{\bits}};
 \node at (\x+1.5,.25) {\value\quad\name};
}
\draw[->] (6,1.45)--(6,1.05);
\end{tikzpicture}
\caption{Byte boundaries and sextet boundaries partition the same bits differently.}\label{base64-bit-split}
\end{figure}

For \texttt{Man}, the decimal bytes are 77, 97 and 110. Substitution gives
\(77\gg2=19\), \(((77\mathbin{\&}3)\ll4)|(97\gg4)=22\),
\(((97\mathbin{\&}15)\ll2)|(110\gg6)=5\), and \(110\mathbin{\&}63=46\).
The standard alphabet maps these values to \texttt{T}, \texttt{W}, \texttt{F} and
\texttt{u}. The order follows the input bit stream, even though the simulator stores
words in little-endian memory. Loading individual bytes avoids reversing a three-byte
group by treating it as a host-endian word.
